\section{Poisson comparison, transport, and logarithmic limits}
\label{sec:log-limits}

\subsection{An independent fragmentation reference and exact transport cost}

For the independent reference, assume the daughter kernel is selfsimilar:
there is a measurable family of measures $B_t$ on $(0,1)$ such that
\begin{equation}
 b_{t,x}=(\theta\mapsto x\theta)_\#B_t,\qquad
 B_t((0,1))=2,\qquad \int_{(0,1)}\theta B_t(\dd\theta)=1.
 \label{eq:selfsimilar-daughters}
\end{equation}
Here $\#$ denotes pushforward. The fraction law $B_t/2$ is independent
of the parent size. Under \eqref{eq:selfsimilar-daughters} the
distribution function of $b_{t,x}/2$ at $z$ is that of $B_t/2$ at
$z/x$, so the quantile of $b_{t,x}/2$ equals $x$ times the quantile
of $B_t/2$. In the construction of \cref{thm:number-process} the
daughter chosen by the mark $u$ is therefore $X_{s-}\Theta$ with
$\Theta$ the quantile of $B_s/2$ at $u$: the retained fraction is a
function of the fragmentation driver's point $(s,u)$ alone, not of
the state. Consequently, on the probability space of
\cref{thm:number-process}, the fragmentation times and fractions form
a marked Poisson random measure $\mathsf P$ of intensity
$2\sigma(t)\dd t\,(B_t/2)(\dd\theta)=\sigma(t)\dd t B_t(\dd\theta)$,
independent of $X_0$ and of the coagulation driver. Define
\begin{equation}
 Z_t=\log X_t,\qquad
 \bar Z_t=Z_0+\int_{(0,t]\times(0,1)}\log\theta\,
                         \mathsf P(\dd s,\dd\theta),\qquad
 \rho_t=\operatorname{Law}(Z_t),\quad
 \bar\rho_t=\operatorname{Law}(\bar Z_t).
 \label{eq:poisson-reference}
\end{equation}
The integral is a finite sum at each finite time, even with infinite
daughter log moments. The increments of $\bar Z_t-Z_0$ are independent;
they are stationary when $\sigma$ is constant and $B_t=B$ is fixed,
as in the setting below. If $\sigma=0$ almost everywhere, the increments
are identically zero regardless of $B_t$. Changes on null sets of times
do not affect the reference law.
This is exactly the number-normalized pure-fragmentation reference,
and no separate solution theory is presupposed for it: with
$\lambda=0$, the linear equation \eqref{eq:weak} is solved directly
by $\bar n_t=N_0e^{F(t)}\operatorname{Law}(e^{\bar Z_t})$, as the
compensator formula for $\mathsf P$ shows, and $\bar n_t$ has count
$N_0e^{F(t)}$ and constant mass $m$ because
$\int\theta B_t(\dd\theta)=1$.

For arbitrary probability measures on $\R$, define the extended
transport cost
\begin{equation}
 \mathsf W_1(\alpha,\beta)
 =\inf_{\gamma\in\Gamma(\alpha,\beta)}
       \int_{\R^2}|z-w|\gamma(\dd z,\dd w)\in[0,\infty],
 \label{eq:extended-W1}
\end{equation}
where $\Gamma$ is the set of couplings. On
$\mathcal P_1(\R)=\{\alpha:\int|z|\alpha(\dd z)<\infty\}$ this
is the ordinary Wasserstein distance. A finite cost between two laws
does not require either law to belong to $\mathcal P_1$. On the
line the classical identity of \citet{Vallender1974},
\[
 \mathsf W_1(\alpha,\beta)=\int_\R|F_\alpha(z)-F_\beta(z)|\dd z
\]
with distribution functions $F_\alpha,F_\beta$, holds in $[0,\infty]$
without moment assumptions, and it is attained by the ordered
(quantile) coupling.

\begin{theorem}[Moment-free exact transport identity]\label{thm:transport}
Under \eqref{eq:selfsimilar-daughters}, the coupling
\eqref{eq:poisson-reference} satisfies $Z_t=\bar Z_t+A_t\geq\bar Z_t$.
Thus $\rho_t$ stochastically dominates $\bar\rho_t$, and
\begin{align}
 \mathsf W_1(\rho_t,\bar\rho_t)
 &=\E A_t=\int_0^t a(s)\dd s\notag\\
 &\leq\frac{H(0)^2}{mN_0}
                     \int_0^t b(s)e^{-\kappa C(s)}\dd s
 \leq\frac{H(0)^2}{\kappa mN_0}.
 \label{eq:transport-exact}
\end{align}
The cost increases to the finite constant $\E A_\infty$.
\end{theorem}

\begin{proof}
The path identity gives both the ordered coupling and the upper bound
on transport cost. Optimality of ordered couplings is classical; we
include a moment-free verification that avoids subtracting infinite
means. Put $c_L(z)=\max(-L,\min(z,L))$. For every coupling $(V,W)$
of $(\rho_t,\bar\rho_t)$, the one-Lipschitz property gives
\[
 \E|V-W|\geq\E c_L(V)-\E c_L(W)
           =\E[c_L(Z_t)-c_L(\bar Z_t)].
\]
For real $z\geq w$, the difference $c_L(z)-c_L(w)$ is the length
of $(w,z]\cap[-L,L]$, so it increases to $z-w$ as $L\to\infty$.
Monotone convergence makes the last expectation tend to $\E A_t$.
This proves the matching lower bound for every coupling, including
ones of infinite cost. The estimates are
\cref{eq:log-rate-bound,eq:controlled-log-tail}.
\end{proof}

\begin{remark}[The boundary at parent dependence]\label{rem:adaptive-marks}
For general $b_{t,x}$ the identity \eqref{eq:path-log-identity} and all
of \cref{thm:controlled-path,cor:controlled-size} remain valid. However,
the realized fractions then depend on $X_{s_j-}$ and hence on its
coagulation history. Removing $A_t$ while keeping those fractions
need not produce a standalone fragmentation process. The independent
Poisson reference and its limits in this section require
\eqref{eq:selfsimilar-daughters}.
\end{remark}

The corresponding bounded-test log equation is useful for identifying
the two mechanisms. For a bounded Borel $\phi$ on $\R$,
\begin{align}
 \pair{\phi}{\rho_t}-\pair{\phi}{\rho_0}
 &=\int_0^t\sigma(s)\int_{\R}\int_{(0,1)}
 [\phi(z+\log\theta)-\phi(z)]B_s(\dd\theta)\rho_s(\dd z)\dd s
 \notag\\
 &\quad+\int_0^t\pair{\phi}{R_s}\dd s,\label{eq:normalized-log}\\
 \pair{\phi}{R_s}
 &:=\frac{\lambda(s)}{N(s)}\iint x
       [\phi(\log(x+y))-\phi(\log x)]n_s(\dd x)n_s(\dd y).
 \label{eq:upward-remainder}
\end{align}
This follows directly from \eqref{eq:number-forward}. The signed
measure $R_s$ has mass zero and variation at most $2b(s)$; it is
nonnegative on increasing tests, and
$|\pair{\phi}{R_s}|\leq\Lip(\phi)a(s)$ for bounded Lipschitz tests.
These are estimates of paired increments. They do not assert separate
absolute log moments for the two positive measures in
\eqref{eq:upward-remainder}.

\subsection{Constant-rate bounds and their critical specialization}

For the remainder of the section, unless stated otherwise, let
$\lambda>0$ and $\sigma\geq0$ be constant, put $b=\lambda m$, and
take $B_t=B$ fixed. Write
\begin{equation}
 r=2\sigma,\qquad \Theta\sim B/2,\qquad Y=\log\Theta<0,\qquad
 \omega=\kappa(b+\sigma),\qquad
 D_0=\frac{\lambda H(0)^2}{\kappa(b+\sigma)N_0}
       \leq\frac b{\kappa(b+\sigma)}.
 \label{eq:constant-log-parameters}
\end{equation}
The symbol $D_0$ here is a transport bound, not relative entropy.
The reference is the compound-Poisson process
$\bar Z_t=Z_0+\sum_{j=1}^{J_t}Y_j$, where $J$ has rate $r$ and
the $Y_j$ are independent copies of $Y$, independent of $J$ and $Z_0$.
The limit $A_\infty$ need not be independent of these variables.

\begin{corollary}[All constant fragmentation rates]\label{cor:constant-cost}
With no logarithmic moment assumptions,
\begin{align}
 \mathsf W_1(\rho_t,\bar\rho_t)&=\E A_t
           \leq D_0(1-e^{-\omega t}),\label{eq:constant-cost}\\
 \E(A_\infty-A_t)&\leq D_0e^{-\omega t}.
 \label{eq:constant-correction-tail}
\end{align}
At criticality $\sigma=b$, count is $N_0$, $\omega=2\kappa b$, and
\begin{equation}
 D_0=\frac{H(0)^2}{2\kappa mN_0}\leq\frac1{2\kappa}.
 \label{eq:critical-cost}
\end{equation}
\end{corollary}

\begin{proof}
In \eqref{eq:log-rate-bound}, $a(t)\leq\lambda H(0)^2
e^{-\omega t}/N_0$. Integrate over $[0,t]$ or $[t,\infty)$.
\end{proof}

The growing one-time cost \eqref{eq:constant-cost} differs from the
decreasing remaining correction \eqref{eq:constant-correction-tail}.
Neither estimate is a bound on raw-size moments or a claim that the
constant $D_0$ is optimal.

\subsection{Laws of large numbers and Gaussian limits}

Pure-fragmentation logarithmic central limit theorems are classical.
For example, \citet{Bertoin2003} studies homogeneous fragmentation
using mass-weighted logarithmic sizes, and
\citet{DoumicEscobedo2016} gives refined critical-fragmentation
asymptotics and a mass-weighted Gaussian limit. Their sampling weights
differ from $B/2$ here. Our number-weighted reference is the
classical many-to-one computation for finite-activity conservative
binary fragmentation, and its limit is an ordinary compound-Poisson
limit; the only new ingredient is the finite nonlinear correction of
\cref{thm:controlled-path}, which transfers that limit to the
additive coagulation--fragmentation equation.

\begin{theorem}[Logarithmic speed and functional central limit theorem]
\label{thm:log-limits}
Assume $\sigma>0$ and fixed $B$.
If $\E|Y|<\infty$, set $\mu=\E Y$. Then
\begin{equation}
 \frac{Z_t}{t}\longrightarrow r\mu
             \quad\text{almost surely}.
 \label{eq:log-SLLN}
\end{equation}
If $\nu_2:=\E Y^2<\infty$, then, for each $S<\infty$,
\begin{equation}
 \left(\frac{Z_{Ts}-r\mu Ts}{\sqrt T}\right)_{0\leq s\leq S}
 \ \Longrightarrow\
 \left(\sqrt{r\nu_2}\,W_s\right)_{0\leq s\leq S}
 \label{eq:functional-CLT}
\end{equation}
in $D([0,S],\R)$ with the Skorokhod $J_1$ topology, where $W$ is
standard Brownian motion. Here $D([0,S],\R)$ is the space of
right-continuous real paths with left limits; $J_1$ permits small
changes in jump times through increasing continuous time changes.
In particular,
\begin{equation}
 \frac{Z_t-r\mu t}{\sqrt t}
          \Longrightarrow\mathcal N(0,r\nu_2).
 \label{eq:weak-CLT}
\end{equation}
No moment assumption on $Z_0$ is needed for these conclusions.
The almost-sure statement \eqref{eq:log-SLLN} and the functional
statement \eqref{eq:functional-CLT} concern the auxiliary process of
\cref{thm:number-process}, one specific coupling of the one-time
laws; the equation itself determines only the one-time limits, such
as \eqref{eq:weak-CLT} and the transport statements below.

If also $\E|Z_0|<\infty$, the one-time convergence
\eqref{eq:weak-CLT} holds in ordinary $\mathsf W_1$, and
\begin{equation}
 \mathsf W_1(\operatorname{Law}(Z_t/t),\delta_{r\mu})
 \leq\frac{\E|Z_0|+D_0}{t}+\sqrt{\frac{r\nu_2}{t}}.
 \label{eq:W1-speed-bound}
\end{equation}
Under only $\E|Y|<\infty$ and $\E|Z_0|<\infty$, ordinary
$\mathsf W_1$ convergence to the speed still holds, without the
displayed square-root rate.
\end{theorem}

\begin{proof}
Put $S_t^Y=\sum_{j\leq J_t}Y_j$. The strong laws for independent
integrable variables and for the Poisson count give
$J_t/t\to r$ and $J_t^{-1}\sum_{j\leq J_t}Y_j\to\mu$ almost
surely. Since $Z_0$ is finite almost surely and $A_t\leq A_\infty$
almost surely, \eqref{eq:path-log-identity} proves \eqref{eq:log-SLLN}.

We give the reduction of the functional assertion to the classical
iid invariance principle, so its moment requirement is explicit.
The increments $\xi_k=S_k^Y-S_{k-1}^Y-r\mu$ are iid with mean zero
and variance $r\nu_2$. Donsker's theorem,
\citet[Theorem~8.2]{Billingsley1999}, gives
$X^N\Rightarrow\sqrt{r\nu_2}\,W$ on $C[0,1]$ as the integer
$N\to\infty$, where $X^N$ is the linear interpolation of
$N^{-1/2}\sum_{k\leq Nu}\xi_k$ between the grid points $u=k/N$; on
$C[0,1]$ the uniform and $J_1$ topologies agree. Rescaling to
$[0,S]$ and passing from integer $N$ to real $T$ are one deterministic
time change. Let $Y_T$ be the linear interpolation of
$T^{-1/2}\sum_{k\leq Ts}\xi_k$ between the grid points $s=k/T$ of
$[0,S]$. With $N=\lceil TS\rceil$,
$Y_T(s)=\sqrt{N/T}\,X^N(\theta_Ts)$, where $\theta_T=T/N\to1/S$ and
$\theta_Ts\leq1$. The map $(x,\theta)\mapsto x(\theta\,\cdot)$ from
$D([0,1],\R)\times(0,1/S]$ to $D([0,S],\R)$ is continuous at every
$(x,\theta)$ with $x$ continuous, because $J_1$ convergence to a
continuous limit is uniform convergence and $x$ is uniformly
continuous; the continuous mapping theorem applied to
$(X^N,\theta_T)$ therefore gives
$Y_T\Rightarrow\sqrt S\sqrt{r\nu_2}\,W(\cdot/S)$ in $D([0,S],\R)$,
and $\sqrt S\,W(\cdot/S)$ is again a standard Brownian motion.
To replace interpolation by the actual compound-Poisson path, let
$V_k=\sum_{k-1<u\leq k}|\Delta S_u^Y|$. Then $V_k$ are iid and
\[
 \E V_k^2=r\nu_2+r^2(\E|Y|)^2<\infty.
\]
On each unit interval, the error between the centered process and
its interpolation is bounded by $2(V_k+r|\mu|)$.
Finite second moment implies
$u^2\Prob(V_1>u)\leq\E[V_1^2\1_{\{V_1>u\}}]\to0$.
For every $\varepsilon>0$, a union bound therefore gives
\[
 \Prob\left(\max_{k\leq\lceil TS\rceil+1}V_k>
                                      \varepsilon\sqrt T\right)
 \leq(\lceil TS\rceil+1)\Prob(V_1>\varepsilon\sqrt T)\to0.
\]
The functional limit follows for the reference process.
Finally,
\begin{equation}
 \sup_{0\leq s\leq S}
 \frac{|Z_{Ts}-S_{Ts}^Y|}{\sqrt T}
 \leq\frac{|Z_0|+A_\infty}{\sqrt T}\to0
 \quad\text{almost surely}.
 \label{eq:uniform-transfer}
\end{equation}
Uniform perturbations tending to zero also tend to zero in $J_1$.
This proves \eqref{eq:functional-CLT}; evaluation at time one gives
\eqref{eq:weak-CLT}.

For the ordinary transport statements, the path identity first gives
finite first log moments at every finite time:
\begin{equation}
 \E|Z_t|\leq\E|Z_0|+rt\E|Y|+\E A_t<\infty.
 \label{eq:log-integrability}
\end{equation}
The centered reference increment has
$\E[(S_t^Y-r\mu t)^2]/t=r\nu_2$. These bounded second moments
give uniform integrability of the first absolute moments of the
scaled reference. Together with its weak Gaussian limit this gives
ordinary $\mathsf W_1$ convergence: by a quantile coupling the variables
converge almost surely to the Gaussian, and uniform integrability
upgrades that convergence to $L^1$. Adding $Z_0/\sqrt t$ costs at
most $\E|Z_0|/\sqrt t$; adding $A_t/\sqrt t$ costs at most
$D_0/\sqrt t$. The triangle inequality proves the asserted ordinary
transport CLT. Cauchy--Schwarz for $S_t^Y-r\mu t$ and the same
coupling give \eqref{eq:W1-speed-bound}.

For the speed with only a first jump moment, truncate the jumps at
$Y^{(L)}=\max(Y,-L)$. The centered truncated sum divided by $t$
tends to zero in $L^1$ by its finite variance. The $L^1$ difference
between the original and truncated centered sums divided by $t$
is at most $2r\E|Y-Y^{(L)}|$, uniformly in $t$. Let first
$t\to\infty$ and then $L\to\infty$. Equation
\eqref{eq:path-log-identity} finishes the proof.
\end{proof}

The Gaussian variance parameter is $r\E Y^2$, the \emph{raw} second
moment of a jump, because the Poisson count fluctuates too. We do not
assert convergence of $\Var(Z_t)/t$: the hypotheses do not even
require a second moment of $Z_0$ or $A_\infty$. Nor does a weak
logarithmic CLT imply convergence to a lognormal density; for
instance, with monodisperse data at size $x_0$ and equal splitting,
the reference stays on the lattice $\log x_0+(\log2)\mathbb Z$ and
the nonlinear size law stays atomic on $\{x_0k2^{-j}:k\geq1,j\geq0\}$.

\begin{corollary}[Moment-free upper speed]\label{cor:infinite-log-speed}
For $\sigma>0$ and fixed $B$, without any log moments,
\[
 \limsup_{t\to\infty}\frac{Z_t}{t}\leq-2\sigma\log2
 \quad\text{almost surely}.
\]
If $\E(-Y)=\infty$, then $Z_t/t\to-\infty$ almost surely.
\end{corollary}

\begin{proof}
For each positive integer $L$, replace every $Y_j$ by
$Y_j^{(L)}=\max(Y_j,-L)\geq Y_j$. The strong law for these bounded
jumps and $A_t\leq A_\infty$ give
$\limsup Z_t/t\leq r\E Y^{(L)}$ on a common probability-one event.
Monotone convergence of the negative parts lets $L\to\infty$.
The resulting bound is $r\E Y$, allowing $-\infty$.
For its universal upper bound, Jensen applied to
$\max(\Theta,e^{-L})$ gives
$\E Y^{(L)}\leq\log\E\max(\Theta,e^{-L})$; the right-hand side
tends to $\log\E\Theta=-\log2$.
\end{proof}

For equal splitting, $\mu=-\log2$ and $\nu_2=(\log2)^2$, so the
speed and Gaussian variance parameter are respectively
$-2\sigma\log2$ and $2\sigma(\log2)^2$. For deterministic
complementary fractions $v$ and $1-v$, they are
$\sigma\log(v(1-v))$ and
$\sigma[(\log v)^2+(\log(1-v))^2]$.

\begin{proposition}[Transfer on every diverging log scale]
\label{prop:general-scale-transfer}
Return to the controlled selfsimilar setting
\eqref{eq:selfsimilar-daughters}. Let $d_t>0$ tend to infinity and
let $c_t$ be any deterministic centering. Then
\begin{equation}
 \E\left|\frac{Z_t-c_t}{d_t}-\frac{\bar Z_t-c_t}{d_t}\right|
 \leq\frac{\E A_\infty}{d_t}\to0.
 \label{eq:general-scale-transfer}
\end{equation}
Every weak limit of the scaled reference laws therefore transfers to
the corresponding scaled number laws. For deterministic paths
$c_T(s)$ and scales $d_T\to\infty$, the same transfer holds in
$D([0,S],\R)$ with $J_1$ whenever the centered reference paths
belong to that space and converge there.
The one-time transfer concerns the laws $\rho_t$, determined by the
equation; the path transfer concerns the auxiliary process of
\cref{thm:number-process}, one specific coupling of the one-time
laws, since the equation itself determines only the one-time limits.
\end{proposition}

\begin{proof}
The one-time bound is $0\leq Z_t-\bar Z_t=A_t\leq A_\infty$.
Markov's inequality and Slutsky's theorem transfer the weak limit.
For paths the uniform difference is at most $A_\infty/d_T$, which
tends to zero almost surely and bounds the $J_1$ distance using the
identity time change.
\end{proof}

This proposition includes a stable limit when the specified reference
satisfies the relevant classical domain-of-attraction assumptions.
It supplies no such domain-of-attraction conclusion for arbitrary
daughter laws, time-dependent marks, or controls.

\subsection{Mean log size, geometric means, and relative entropy}

\begin{proposition}[Exact logarithmic identities]\label{prop:log-means}
In the constant-rate, fixed-$B$ setting assume
$\E|Z_0|<\infty$ and, when $\sigma>0$, $\E|Y|<\infty$.
Write $\ell=r\E Y$ when $\sigma>0$, and $\ell=0$ when
$\sigma=0$. Then
\begin{equation}
 \E Z_t=\E Z_0+\ell t+\E A_t.
 \label{eq:mean-log}
\end{equation}
The geometric mean $G(t):=\exp(\int\log x\,\eta_t(\dd x))$
therefore satisfies
\begin{equation}
 e^{-\ell t}G(t)\longrightarrow
 G_*:=\exp(\E Z_0+\E A_\infty)\in(0,\infty),\qquad
 0\leq\log G_*-\log(e^{-\ell t}G(t))\leq D_0e^{-\omega t}.
 \label{eq:geometric-mean}
\end{equation}
The relative entropy of number against mass sampling is finite and
obeys
\begin{align}
 D_{\rm KL}(\eta_t\Vert\pi_t)
 &:=\int_E\log\left(\frac{\dd\eta_t}{\dd\pi_t}\right)\dd\eta_t
 =\log\frac m{N(t)}-\E Z_t\notag\\
 &=D_{\rm KL}(\eta_0\Vert\pi_0)
       +(b-\sigma-\ell)t-\E A_t.
 \label{eq:KL-identity}
\end{align}
For $\sigma>0$, $\mu=\E Y\leq-\log2$, so its asymptotic linear
growth rate is
$b-\sigma-r\mu\geq b+\sigma(2\log2-1)>0$.
At criticality this rate is $-2b\mu$.
\end{proposition}

\begin{proof}
Integrability follows from \eqref{eq:log-integrability}, with the jump
sum absent if $\sigma=0$. Taking expectations in the finite-cost
path identity proves \eqref{eq:mean-log}; the correction-tail bound
proves \eqref{eq:geometric-mean}. The sampling laws are mutually
absolutely continuous, with
$\dd\eta_t/\dd\pi_t=m/(N(t)x)$. Its logarithm is integrable by
the established first log moment, so integration and
\cref{eq:count,eq:mean-log} give \eqref{eq:KL-identity}.
Finally Jensen's inequality gives
$\mu\leq\log\E\Theta=-\log2$.
\end{proof}

No mass-weighted logarithmic moment such as
$\int x|\log x|n_0(\dd x)$ is needed here. Without the stated first
log moments we retain the paired transport identity, but make no
finite mean-log, positive geometric-mean-prefactor, or finite-entropy
claim. The entropy identity expresses the standard change of sampling
weight; it is not an additional entropy-production estimate.

\subsection{The pure additive endpoint and the classical Borel solution}

\begin{corollary}[Pure coagulation]\label{cor:pure-coag}
Let $\sigma=0$ and $\lambda>0$ be constant. Then $Z_t$ increases to
a finite real $Z_\infty$ almost surely and is eventually constant.
Writing $\rho_\infty=\operatorname{Law}(Z_\infty)$, for $0\leq s\leq t$,
\begin{align}
 \mathsf W_1(\rho_s,\rho_t)&=\int_s^t a(u)\dd u,\label{eq:pure-cost}\\
 \mathsf W_1(\rho_t,\rho_\infty)&=\int_t^\infty a(u)\dd u
                   \leq D_0e^{-\kappa bt}.
 \label{eq:pure-limit}
\end{align}
These statements require no initial log moment. In particular the
number-normalized size laws have a proper weak limit on $E$.
The transport identities and this weak limit are properties of the
equation; the almost-sure monotone convergence concerns the
constructed auxiliary process.
\end{corollary}

\begin{proof}
Now $Z_t=Z_0+A_t$, and \cref{thm:controlled-path} supplies the
finite limit and finite total jumps. Apply the clipped-test
optimality argument of \cref{thm:transport} to the ordered pairs
$(Z_t,Z_s)$ and $(Z_\infty,Z_t)$. Their expected differences are
the displayed integrals. Almost-sure convergence and continuity of
the exponential give the size-law convergence.
\end{proof}

The existence of a proper number-normalized limit at this endpoint
has classical precedents. For $n_0=\delta_1$, $m=1$, and $\lambda=1$,
the monodisperse additive solution is
\begin{equation}
 n_t(\{k\})=e^{-t}\frac{[k(1-e^{-t})]^{k-1}}{k!}
                        e^{-k(1-e^{-t})},\qquad k\geq1,\ t>0.
 \label{eq:classical-Borel}
\end{equation}
This is the classical Golovin solution as recorded in
\citet[equation~(3)]{Bertoin2009}; branching representations of the
additive solution are developed in \citet{DeaconuTanre2000}.
Its number law is $\operatorname{Borel}(1-e^{-t})$ and tends to
$\operatorname{Borel}(1)$, the proper law with probabilities
$p_k=e^{-k}k^{k-1}/k!$. One can see that this critical law is proper
as the total progeny of a Poisson-one Galton--Watson tree: its
extinction probability is the smallest fixed point of
$q=e^{q-1}$ in $[0,1]$, namely one. Stirling's formula gives
$p_k\sim(2\pi)^{-1/2}k^{-3/2}$, so its mean is infinite even
though its first logarithmic moment is finite. At finite times the
number mean size is $m/N(t)=e^t$. This classical example shows why
the finite log transport limit does not control the arithmetic
mean, and why the pure-coagulation endpoint itself is not a new
limiting phenomenon.
